\section{Methods}
\label{sec:methods}

\subsection{Conserved customer-to-geography fuel ledger}

Every OCAC record retains its period, customer category, fuel type, reported tonnes, source
file and revision status. It is assigned across six non-overlapping destinations: IPCC
domestic civil aviation ($d$), international bunkers ($i$), military/government 1A5b ($m$),
general aviation ($g$), export/other ($o$) and unresolved residual ($u$). For each source row
$r$,
\begin{equation}
F_r=F_{r,d}+F_{r,i}+F_{r,m}+F_{r,g}+F_{r,o}+F_{r,u}.
\label{eq:rowclosure}
\end{equation}
Closure is asserted before aggregation. Missing evidence increases $F_{r,u}$; it never triggers
a national rescaling.

Three structural mappings test the customer-to-geography boundary. A literal mapping assigns
resident-airline JP-1 to domestic aviation. A movement-share mapping divides it according to
domestic and international flight stages. The central calibrated mapping uses the OECD
resident-carrier share
\begin{equation}
s_d=\frac{E^{\rm OECD}_{\rm resident,domestic}}
{E^{\rm OECD}_{\rm resident,domestic}+E^{\rm OECD}_{\rm resident,intl.from}},
\label{eq:share}
\end{equation}
which is 25.674\% in FY2023--24. Foreign-customer JP-1 is assigned principally to international
bunkers; government aviation JP-1 and JP-8 are military/government candidates; 100LL is a
general-aviation candidate; and government-other JP-1 remains unresolved. Alternative
carrier-specific movement and modelled-fuel shares preserve separate PIA and other-airline
networks while remaining constrained to the same aggregate OECD share. These mappings are
reported as scenarios, not pooled into a confidence interval.

Three definitions are fixed as accounting invariants. PCAA movements are divided by two to
obtain domestic stages because each flight generates a departure and an arrival. Domestic
status requires both endpoints in Pakistan, irrespective of customer or carrier residence.
Domestic greenhouse gases use full-flight fuel; international airport LTO emissions remain a
separate local-air-quality product \citep{ipcc2006energy}.

\subsection{Domestic activity and fuel models}

The aggregate movement model (M0) multiplies domestic stages by a representative distance and
year-varying fleet intensity:
\begin{equation}
E_y=N_y\,\bar d_y\,I_y,
\end{equation}
where $N_y$ is flight stages, $\bar d_y$ mean distance and $I_y$ fuel and emission intensity.
The route-resolved model (M1) replaces the representative stage with an annual origin--
destination matrix constrained to PCAA airport margins, airline departures and routes
documented for FY2023--24. Unsupported movements remain an unallocated residual. Minimum and
maximum feasible distance are solved over the same route graph and control margins.

The route--aircraft model (M2) assigns probabilistic aircraft-family shares by operator and
route and generates phase-resolved synthetic profiles with OpenAP \citep{sun2020openap}.
Operator passengers, cargo and mail constrain central payload; route detour, taxi time and
aircraft assignment remain uncertain. Unsupported aircraft use a separately identified ICAO
factor, while route, operator and endpoint residuals are completed through explicit lower-tier
methods. A common-support comparison applies all three models to the same 17,491 stages before
those residuals are added, preventing model resolution from being confused with coverage.
Synthetic route strata are never labelled as observed flights.

The flight-stage ledger is source-neutral: observed records may replace reconstructed strata
without changing accounting or emissions code. The detailed route optimiser, source tiers,
OpenAP range checks, ADS-B endpoint rules, payload construction and coverage activation gate
are provided in the supplement.

\subsection{Evaluation and comparison design}

OCAC cumulative monthly values are first-differenced and required to reproduce each final
annual source row. OECD records are aggregated to identical July--June windows. For component
$j$ and month or year $t$,
\begin{equation}
R_{t,j}=\frac{F^{\rm OCAC}_{t,j}\,3.16}{E^{\rm OECD}_{t,j}}.
\end{equation}
We compare resident, foreign and combined commercial accounts and calculate the fraction of
opposite component differences that cancels in their total:
\begin{equation}
C=1-\frac{|\Delta_{\rm resident}+\Delta_{\rm foreign}|}
{|\Delta_{\rm resident}|+|\Delta_{\rm foreign}|}.
\end{equation}
Five complete fiscal years support annual comparison and 42 matched months test whether annual
closure is temporally stable.

The central fuel bridge uses OECD's resident share and therefore tests consistency, not
independent accuracy. The strongest gross closure compares OCAC commercial JP-1 with OECD
combined territory aviation. Aggregate movement, route-resolved and route--aircraft models are
compared externally with OECD domestic territory emissions. Pakistan's CRT tests the official
boundary implementation; Climate TRACE tests fiscal-year magnitude and airport ordering but
may share surveillance or performance inputs. PCAA movement, passenger and route constraints
target specific internal failure modes.

Public ADSB.lol and ERA5 experiments remain supporting diagnostics. The former tests whether a
small public trace sample can replace reconstructed stages; the latter tests the sensitivity of
representative OpenAP flight duration to route winds. Neither is expanded into an observed
annual inventory or counted as fuel validation.

\subsection{Parametric and structural uncertainty}

Parametric and structural uncertainty are evaluated separately. The aggregate movement model
retains its legacy activity/fleet representation term. The completed route--aircraft model uses
20,000 flight-stage Monte Carlo draws within each structural configuration. Aircraft family is
shared within operator--route strata, detour across reciprocal routes, taxi time by airport and
carbon conversion globally. Payload is sampled around PCAA passenger, cargo and mail controls.
CH4, N2O and net calorific value enter only the secondary \COe{} result.

Alternative feasible route matrices and unsupported-activity completion methods define a
structural envelope; customer-to-geography mappings remain a separate accounting scenario.
The route--aircraft model does not inherit the aggregate model's $\pm60$\% representation term.
A drop-one variance screen identifies which named variables drive its interval. Automated QA/QC
asserts source-row and monthly closure, route-bound ordering, common support, coverage
completion, phase-fuel closure, movement balance and agreement between generated outputs and
the manuscript. The complete uncertainty register and assertion catalogue are supplementary.
